%% kapa: exam=1 round=29 year=2018 period=2 form=A subject=accounting subject_ko=회계학 seq=1 number=41
%% source: Q-Net 감정평가사 기출문제 (cst00309) artlSeq=5202808 file=extracted/1008/5202808/2018년 제29회 감정평가사 1차시험 기출문제/2018년 제29회 감정평가사 1차시험 기출문제/2교시 A형.pdf
%% conversion: kordoc 4.24.51 -> tools/gen1.py (auto); review=proofread-sonnet; flags=none
\begin{question}{41}
재무보고를 위한 개념체계에서 유용한 재무정보의 질적 특성에 관한 설명으로 옳은 것은?
\begin{choices}
  \item 재무정보가 예측가치를 갖기 위해서 그 자체가 예측치 또는 예상치일 필요는 없다.
  \item 계량화된 정보가 검증가능하기 위해서 단일 점추정치이어야 한다.
  \item 완벽하게 표현충실성을 위해서는 서술은 완전하고, 검증가능하며, 오류가 없어야 한다.
  \item 재무정보에 예측가치가 있다면 그 재무정보는 나타내고자 하는 현상을 충실하게 표현한다.
  \item 재고자산평가손실의 인식은 보수주의 원칙이 적용된 것이며, 보수주의는 표현충실성의 한 측면으로 포함할 수 있다.
\end{choices}
\end{question}
\answer{1}
